\section{Geometric flagships}
\label{sec:flagships}

\subsection{The twisted-cubic--axis family}

Let $q=3^h$ and work in $\mathrm{PG}(3,q)$.  Put
\[
 T_q=\{C(t)=(1:t:t^2:t^3):t\in\F_q\},
\]
\[
 L_q=\{A(u)=(0:1:u:0):u\in\F_q\}
 \cup\{A_\infty=(0:0:1:0)\}.
\]
The line $L_q$ is the characteristic-three axis of the osculating-plane
pencil \cite{DavydovEtAl2021}; all statements below follow directly from
the displayed coordinates.

\begin{theorem}[Cubic--axis port]
\label{thm:cubic}
The row code on $T_q\sqcup L_q$ is $[2q+1,4,q-1]_q$.  Its minimal circuits
of size at most four are the triples on $L_q$ and the sets
\[
 \{C(s),C(t),C(u),(0:e_1:e_2:0)\}
\]
for distinct $s,t,u$, where $e_1=s+t+u$ and
$e_2=st+su+tu$.  Every cubic coordinate has
\[
 (\nu,\tau)=\left(\frac{q-1}{2},q-2\right).
\]
If $Z_3(q)$ is the maximum size of a three-term-zero-sum-free subset of
$(\F_q,+)$, then
\[
 (\nu,\tau)_{A_\infty}
 =\left(\frac{5q-3}{6},2q-1-Z_3(q)\right),
\]
\[
 (\nu,\tau)_{A(a)}
 =\left(\frac{5q-9}{6},2q-2-Z_3(q)\right).
\]
In particular every coordinate satisfies $\tau>\nu$ for $q\geq9$.
\end{theorem}

\begin{proof}
A plane containing $L_q$ contains at most one cubic point by Frobenius
injectivity; a plane not containing it meets the line once and the cubic at
most three times.  This gives the code parameters.  Vandermonde determinants
classify the displayed small circuits.  At an axis target the repair clutter
is a disjoint union of $K_q$ and a zero-sum-triple hypergraph, so matching and
transversal numbers add; complements of transversals are precisely
zero-sum-free sets.  At a cubic target, inversion relabels the other cubic
helpers by $\F_q^*$.  A multiplicative-generator pairing gives $(q-1)/2$
repairs with distinct axis completions; no larger matching is possible because
each repair uses two of the $q-1$ cubic helpers.  The shifted-inverse
complement lemma bounds a repair-free helper set by $q+2$, attained by all
axis helpers together with one cubic helper.  Taking complements gives
$\tau=q-2$.  The stated strict inequalities follow from the elementary bound
$Z_3(q)\leq q-2$ and the displayed formulas.
\end{proof}

Restoring the cubic point at infinity gives a
$[2q+2,4,q]_q$ code whose full minimal port is visible at radius four and has
uniform rows
\[
 \left(\frac{q-1}{2},q-1\right)\text{ on the cubic},
 \qquad
 \left(\frac{5q-3}{6},2q-3\right)\text{ on the axis}.
\]
This is the inner code in the strict weighted-transfer example above.
The exact $q=9$ rows are recorded in Appendix~\ref{sec:finite-evidence}.

\subsection{The quartic--nucleus harmonic family}

Let
\[
 V(t)=(1,t,t^2,t^3,t^4),\qquad V(\infty)=e_4,
 \qquad N=e_2,
\]
and let $S_q=\{V(t):t\in\mathbb P^1(\F_q)\}\cup\{N\}$.
The classical normal-rational-curve nucleus formula identifies $N$ as the
hyperplane nucleus in characteristic three \cite[Theorem~1]{GmainerHavlicek2013}.

\begin{theorem}[Quartic--nucleus port]
\label{thm:harmonic}
For $q=3^h\geq9$, the row code $Q_q$ on $S_q$ has parameters
\[
 [q+2,5,q-3]_q,\qquad d(Q_q^\perp)=5.
\]
Its minimal circuits of size at most five are exactly
\[
 \{N\}\cup B,
\]
where $B$ is a harmonic quadruple of $\mathbb P^1(\F_q)$.  These quadruples
form a Steiner system $S(3,4,q+1)$, the characteristic-three four-set orbit
in the standard $\operatorname{PGL}(2,q)$ design construction
\cite[Section~2]{Tricot2024}.  Hence the radius-four repairs are all
$B$ at the nucleus and all $\{N\}\cup(B\setminus\{x\})$ at a curve target
$x\in B$.  The nucleus has
\[
 \frac{(q+1)q(q-1)}{24}
\]
radius-four repairs, and every curve target has $q(q-1)/6$.  Every curve
target has $(\nu,\tau)=(1,1)$.
\end{theorem}

\begin{proof}
For distinct finite $a,b,c,d$, the hyperplane through their quartic columns
contains $N$ exactly when $e_2(a,b,c,d)=0$; with $\infty$ present the condition
is $a+b+c=0$.  Given three finite points, if $e_1=a+b+c\neq0$, the unique
completion is $d=-e_2/e_1$; if $e_1=0$, the unique completion is $\infty$.
In the first case, substituting $d=a$, for example, changes the harmonic
equation to $(a-b)(a-c)=0$, and similarly for $b$ and $c$; hence the
completion is distinct.  In the second case, writing $c=-a-b$ gives
$e_2(a,b,c)=-(a-b)^2\neq0$, so no finite completion exists.  A triple
$\{\infty,a,b\}$ has the unique completion $-a-b$, distinct from $a,b$
in characteristic three.  This proves the Steiner property.  Since
$\{\infty,0,1,-1\}$ is one block and $\operatorname{PGL}(2,q)$ is
3-transitive, unique completion also identifies the blocks with its
projective orbit.

Any five curve points are independent by the extended Vandermonde
determinant.  Expanding the determinant of $N,V(a),V(b),V(c),V(d)$ along
the nucleus row gives
\[
 \Delta(a,b,c,d)e_2(a,b,c,d),
\]
and with infinity present it gives $\Delta(a,b,c)e_1(a,b,c)$.  The
Vandermonde factor is nonzero for distinct parameters.  Thus
$\{N\}\cup B$ is dependent exactly for a harmonic block.  A complementary
minor shows that $N$ with any three distinct curve points is independent;
in the zero-sum finite branch its determinant reduces to
$-\Delta(a,b,c)e_2(a,b,c)$, which is nonzero by the calculation above.
Deleting any point therefore restores independence.  These are all circuits
of size at most five, so the dual distance is five.

A nonzero hyperplane restricts to a nonzero polynomial of degree at most four
on the normal rational curve, counting infinity as a root when the quartic
coefficient vanishes.  It therefore contains at most four curve points.  If
it also contains $N$, equality occurs exactly on a harmonic block, so the
largest section of $S_q$ has five points.  Hence the primal distance is
$(q+2)-5=q-3$.  Finally, the Steiner count is
$\binom{q+1}{3}/\binom43$, and the number of blocks through a fixed curve
point is $\binom q2/\binom32$.  The circuit--repair correspondence gives the
stated ports and counts.
\end{proof}

The inner-dual half of the coarse transfer gate is automatic at radius four
because $5<2d(Q_q^\perp)=10$.  Theorem~\ref{thm:prescribed} therefore places the
entire harmonic port with positive density in asymptotically good
fixed-$\F_q$ families.

\subsection{Reliability, EXIT, and the geometric contrast}

The nucleus repairs are the blocks of the Steiner system
\(S(3,4,q+1)\).  At a curve target, every radius-four repair contains the
nucleus.  Thus the nucleus port is a parallel family of four-helper repairs,
whereas the curve port has a compulsory series helper.  The reliability and
bounded-EXIT theorem applies to both ports without an exact finite profile.
Exact field-nine profiles are computational refinements recorded in
Appendix~\ref{sec:finite-evidence}; they are not used in the body theorem chain.

Finally, the harmonic small-circuit closure has a deterministic nucleus gate.
Let $D(S)$ close curve set $S$ under ``three points of a harmonic block add
the fourth.''  Then
\[
 \operatorname{cl}_4(S)=S\quad\text{if $S$ contains no block},
\qquad
 \operatorname{cl}_4(S)=\{N\}\cup D(S)\quad\text{otherwise},
\]
and $\operatorname{cl}_4(S\cup\{N\})=\{N\}\cup D(S)$.
Indeed, without $N$ no curve point can be repaired, while a complete block
$B\subseteq S$ repairs $N$.  Once $N$ is present, the curve-target repairs
are exactly $N$ together with three points of a harmonic block, so iteration
is precisely harmonic completion.  This proves all three identities.

Appendix~\ref{sec:finite-evidence} gives a field-nine block-free rank-five
spanning set and its one-round completion after adjoining $N$.  These witnesses
illustrate the nucleus gate but are not inputs to the closure identities, an
asymptotic propagation claim, or a threshold theorem.
